% Extracción quirúrgica del propietario V2 05e: líneas 325--447.
% El resto permanece en este archivo como procedencia, pero no se publica.
\iffalse
\chapter[Macrostructure déterminantielle interne]{Macrostructure déterminantielle interne engendrée par le continu discret}
\label{chap:detint-macroestructura}

\begin{thesisbox}
La branche \({\rm det}\) part exclusivement d'histoires survivantes de
\(\mathcal C_{\rm cont}^{\rm disc}\). Le nerf des extensions, la diagonale
Alexander--Whitney, le complexe local, l'unité, les deux publications, le
filler, la réduction terminale et la tour \(3\)-adique se construisent avant
toute évaluation ultérieure. Les conditions de bilan et d'acyclicité sont
incorporées au domaine typé et ne sont jamais tacitement présupposées.
\end{thesisbox}

\section{Domaine déterminantiel survivant}

Soit
\begin{equation}\label{eq:detint-an}
 A_n:=\mathbf Z/3^n\mathbf Z,
 \qquad
 \rho_{n+1,n}:A_{n+1}\longrightarrow A_n.
\end{equation}
Pour \(x\in\mathcal C_{\rm cont}^{\rm disc}\), soit
\(\mathsf T_m(x)\) la catégorie finie des préfixes TPK de profondeur \(m\)
compatibles avec \(x\). Chaque objet conserve feuillet, orientation, résidu,
quotient, retenue entière, mémoire, frontière, chemin, multisection, survie
et terminal ; ses morphismes sont les extensions admissibles qui conservent ces
coordonnées.

De chaque objet \(y\in\mathsf T_m(x)\), on extrait les entiers internes
\begin{equation}\label{eq:detint-c-v-lifts}
 \widetilde c_y,\widetilde v_y\in\mathbf Z,
 \qquad
 c_{y,n}:=\widetilde c_y\bmod3^n,
 \qquad
 v_{y,n}:=\widetilde v_y\bmod3^n.
\end{equation}
La première coordonnée est la retenue non orientée et la seconde, le coefficient
orienté de frontière. Le renversement satisfait
\begin{equation}\label{eq:detint-or-cv}
 c_{\bar y,n}=c_{y,n},
 \qquad
 v_{\bar y,n}=-v_{y,n}.
\end{equation}

On écrit
\begin{equation}\label{eq:detint-z3-k3}
 \mathbf Z_3:=\varprojlim_n A_n,
 \qquad
 K_3:=\operatorname{Frac}(\mathbf Z_3).
\end{equation}
La restriction déterminantielle est définie par
\begin{equation}\label{eq:detint-dominio}
\begin{split}
 \mathcal C_{\rm cont}^{\rm det}:=
 \bigl\{x\in\mathcal C_{\rm cont}^{\rm disc}:{}&
 x\text{ survit à toute profondeur};\\
 &\widehat P_x^{\rm red}\text{ est fini et basé};\\
 &\chi(\widehat P_x^{\rm red})=0;\\
 &H_*(\widehat P_x^{\rm red}\otimes_{\mathbf Z_3}K_3)=0;\\
 &\text{la réduction canonique produit un bloc carré}\\
 &\text{stable par raffinement}\bigr\}.
\end{split}
\end{equation}
Les objets \(\widehat P_x^{\rm red}\) et le bloc carré seront construits
dans \cref{sec:detint-terminal}. Par conséquent,
\eqref{eq:detint-dominio} est une condition d'appartenance vérifiable et non
une conclusion insérée dans la définition des applications.

\section{Nerf, Alexander--Whitney et counité}

On définit le complexe normalisé
\begin{equation}\label{eq:detint-gmn}
 G_{m,n}(x):=
 N_*^{\rm norm}\!\left(N\mathsf T_m(x);A_n\right).
\end{equation}
Pour un simplexe non dégénéré
\(\sigma=[y_0\to y_1\to\cdots\to y_q]\), la diagonale est
\begin{equation}\label{eq:detint-aw}
 \Delta_{\rm AW}(\sigma)
 =\sum_{j=0}^q
 [y_0\to\cdots\to y_j]\otimes
 [y_j\to\cdots\to y_q].
\end{equation}
La counité est donnée par
\begin{equation}\label{eq:detint-counidad-g}
 \epsilon_G([y])=1,
 \qquad
 \epsilon_G(\sigma)=0\quad(q>0).
\end{equation}
En particulier,
\begin{equation}\label{eq:detint-vertice-grouplike}
 \Delta_{\rm AW}[y]=[y]\otimes[y],
 \qquad
 \epsilon_G([y])=1.
\end{equation}

Si \(m'\ge m\) et \(n'\ge n\), la troncature et la réduction des coefficients
induisent
\begin{equation}\label{eq:detint-rmn}
 R_{m',m}^{n',n}:G_{m',n'}(x)\longrightarrow G_{m,n}(x),
\end{equation}
et l'on vérifie
\begin{equation}\label{eq:detint-aw-natural}
 R\partial=\partial R,
 \qquad
 (R\otimes R)\Delta_{\rm AW}=\Delta_{\rm AW}R,
 \qquad
 \epsilon_GR=\rho_{n',n}\epsilon_G.
\end{equation}
Le chemin ne se réduit pas à ses extrémités : le simplexe entier reste dans
\(G_{m,n}(x)\).

\section{Complexe local interne}

Pour \(c\in A_n\), on définit
\begin{equation}\label{eq:detint-p0-grados}
 P^0_{n,1}(c)=A_nf,
 \qquad
 P^0_{n,0}(c)=A_nr\oplus A_ne\oplus A_nb,
 \qquad
 P^0_{n,-1}(c)=A_ng,
\end{equation}
avec différentielle
\begin{equation}\label{eq:detint-p0-diferencial}
 \partial f=3ce+b,
 \qquad
 \partial r=0,
 \qquad
 \partial e=g,
 \qquad
 \partial b=-3cg.
\end{equation}
L'identité de chaîne est littérale :
\begin{equation}\label{eq:detint-d2}
 \partial^2f=3c\,\partial e+\partial b=3cg-3cg=0,
 \qquad
 \partial^2e=\partial^2b=\partial^2r=0.
\end{equation}

L'augmentation
\begin{equation}\label{eq:detint-epsilon-p}
 \epsilon_P:P_n^0(c)\longrightarrow A_n[0]
\end{equation}
est fixée par
\[
 \epsilon_P(r)=1,
 \qquad
 \epsilon_P(e)=\epsilon_P(b)=\epsilon_P(f)=\epsilon_P(g)=0.
\]
Le complexe possède la coalgèbre différentielle
\begin{equation}\label{eq:detint-coalgebra-p}
 \Delta_P(r)=r\otimes r,
 \qquad
 \Delta_P(q)=q\otimes r+r\otimes q
 \quad(q\in\{e,b,f,g\}).
\end{equation}
Ainsi, \(r\) est group-like et les quatre autres générateurs sont primitifs
relativement à \(r\). Les équations
\eqref{eq:detint-p0-diferencial} prouvent que \(\Delta_P\) commute avec la
différentielle tensorielle.

\section{Système local d'extensions et ledger de mémoire}

Pour une extension \(\alpha:y\to y'\), nous définissons
\begin{equation}\label{eq:detint-psi}
 \Psi_\alpha:P_n^0(c_y)\longrightarrow P_n^0(c_{y'})
\end{equation}
par
\begin{equation}\label{eq:detint-psi-generadores}
\begin{gathered}
 \Psi_\alpha(r_y)=r_{y'},\quad
 \Psi_\alpha(e_y)=e_{y'},\quad
 \Psi_\alpha(g_y)=g_{y'},\quad
 \Psi_\alpha(f_y)=f_{y'},\\
 \Psi_\alpha(b_y)=b_{y'}+3(c_{y'}-c_y)e_{y'}.
\end{gathered}
\end{equation}
C'est une application de complexes parce que
\begin{equation}\label{eq:detint-psi-b-chain}
 \partial\Psi_\alpha(b_y)
 =-3c_{y'}g+3(c_{y'}-c_y)g
 =-3c_yg
 =\Psi_\alpha(\partial b_y),
\end{equation}
et
\begin{equation}\label{eq:detint-psi-f-chain}
 \Psi_\alpha(\partial f_y)
 =3c_ye+b+3(c_{y'}-c_y)e
 =3c_{y'}e+b
 =\partial\Psi_\alpha(f_y).
\end{equation}
De plus,
\begin{equation}\label{eq:detint-psi-functorial}
 \Psi_\beta\Psi_\alpha=\Psi_{\beta\alpha}.
\end{equation}

Le changement de la coordonnée orientée n'est pas effacé. Il est enregistré par
\begin{equation}\label{eq:detint-k-alpha}
 K_\alpha:=(v_{y'}-v_y)f_{y'}.
\end{equation}
Pour des extensions composables,
\begin{equation}\label{eq:detint-k-coherencia}
 K_{\beta\alpha}=K_\beta+\Psi_\beta K_\alpha.
\end{equation}
Cette égalité est le ledger exact de l'accroissement de frontière et rend visible
la mémoire qu'une composition purement nominale perdrait.

Le diagramme local complet est la construction de Grothendieck
\begin{equation}\label{eq:detint-grothendieck-local}
 \int_{y\in\mathsf T_m(x)}P_n^0(c_y).
\end{equation}
Dans celle-ci, une flèche au-dessus de \(\alpha:y\to y'\) porte comme composante linéaire
\(\Psi_\alpha\) et comme cohérence \(K_\alpha\). Lorsqu'elle est appliquée à un
élément supporté en un sommet, on utilise la notation
\begin{equation}\label{eq:detint-psi-soporte-vertice}
 \Psi_\alpha(p,[y]):=(\Psi_\alpha p,[y']).
\end{equation}
L'arête \([y\to y']\) reste dans la coordonnée de chemin du nerf. Cette
convention sera celle utilisée dans les formules de naturalité de
\(z_\pm\) et \(h_*\) ; on n'affirme pas une application globale qui efface les autres
simplexes de \(G_{m,n}(x)\).

\section{Pullback, unité, racine, publications et filler}

Le pullback de complexes est
\begin{equation}\label{eq:detint-pgen}
 P^{\rm gen}_{y;m,n}
 :=P_n^0(c_y)\times_{A_n[0]}G_{m,n}(x)
 =\{(p,\gamma):\epsilon_P(p)=\epsilon_G(\gamma)\}.
\end{equation}
Sa différentielle est
\begin{equation}\label{eq:detint-pgen-d}
 \partial(p,\gamma)=(\partial p,\partial\gamma).
\end{equation}
L'unité commune associée au sommet \(y\) est
\begin{equation}\label{eq:detint-unidad}
 \mathbf1_y:=(r,[y]).
\end{equation}
Ses deux projections sont group-like par
\eqref{eq:detint-vertice-grouplike} et
\eqref{eq:detint-coalgebra-p} ; c'est le sens précis d'unité
group-like compatible du pullback.

La projection racine est l'application de complexes
\begin{equation}\label{eq:detint-proyeccion-root}
 \varpi_{\rm root}:P^{\rm gen}_{y;m,n}\longrightarrow
 \operatorname{Root}_{y;m,n},
 \qquad
 \varpi_{\rm root}(p,\gamma):=(\epsilon_P(p)r,\gamma).
\end{equation}
On définit les deux publications internes
\begin{equation}\label{eq:detint-zpm}
 z_-(y):=(r,[y]),
 \qquad
 z_+(y):=(r+3c_yv_ye+v_yb,[y]).
\end{equation}
Toutes deux sont des cycles :
\begin{equation}\label{eq:detint-zpm-ciclos}
 \partial z_-(y)=0,
 \qquad
 \partial z_+(y)=3c_yv_yg-3c_yv_yg=0.
\end{equation}
Leurs racines coïncident :
\begin{equation}\label{eq:detint-raiz-comun}
 \varpi_{\rm root}z_-(y)
 =(r,[y])
 =\varpi_{\rm root}z_+(y).
\end{equation}

Le filler se construit directement :
\begin{equation}\label{eq:detint-hstar}
 h_*(y):=(-v_yf,0)\in P^{\rm gen}_{y;m,n,1}.
\end{equation}
Alors
\begin{equation}\label{eq:detint-filler-identidad}
 \partial h_*(y)
 =(-3c_yv_ye-v_yb,0)
 =z_-(y)-z_+(y).
\end{equation}
En particulier,
\begin{equation}\label{eq:detint-clases-iguales}
 [z_-(y)]=[z_+(y)]
 \quad\text{dans}\quad H_0(P^{\rm gen}_{y;m,n}),
\end{equation}
et l'égalité conserve sa chaîne témoin.

Sous une extension \(\alpha:y\to y'\), on a
\begin{equation}\label{eq:detint-z-natural}
 z_-(y')=\Psi_\alpha z_-(y),
 \qquad
 z_+(y')-\Psi_\alpha z_+(y)=\partial K_\alpha,
\end{equation}
et
\begin{equation}\label{eq:detint-h-natural}
 h_*(y')-\Psi_\alpha h_*(y)=-K_\alpha.
\end{equation}
Par conséquent, la naturalité n'identifie pas des valeurs distinctes de \(v\) : elle enregistre
leur différence au moyen de l'homotopie cohérente \(K_\alpha\).

\section{Orientation interne}

Sur \(P_n^0(c)\), on définit
\begin{equation}\label{eq:detint-omega-p}
 \Omega(r)=r,
 \qquad
 \Omega(q)=-q\quad(q\in\{e,b,f,g\}).
\end{equation}
C'est un automorphisme de complexes et de la coalgèbre relative. Dans le nerf,
l'inversion orientée est
\begin{equation}\label{eq:detint-omega-g}
 \mathfrak o_G[y_0\to\cdots\to y_q]
 :=(-1)^{q(q+1)/2}
 [\bar y_q\to\cdots\to\bar y_0].
\end{equation}
Avec \eqref{eq:detint-or-cv}, ces actions donnent
\begin{equation}\label{eq:detint-or-z-h}
 \Omega z_-(y)=z_-(\bar y),
 \qquad
 \Omega z_+(c_y,v_y)=z_+(c_{\bar y},v_{\bar y}),
 \qquad
 \Omega h_*(y)=h_*(\bar y).
\end{equation}
L'orientation agit sur le chemin, le complexe, les publications, le filler et la droite
déterminantielle ; ce n'est pas une étiquette sans action.

\fi
\section{Relèvement du conoyau au moyen de la tour de Bockstein et conservation de la retenue}
\label{sec:detint-terminal}

Seule la racine commune est supprimée :
\begin{equation}\label{eq:detint-p-reducido}
 \overline P_{y;m,n}
 :=P^{\rm gen}_{y;m,n}/A_n\mathbf1_y,
 \qquad
 \widehat P_{y;m}^{\rm red}:=\varprojlim_n\overline P_{y;m,n}.
\end{equation}
La base est ordonnée d'abord selon l'ordre TPK des chemins, puis selon
\begin{equation}\label{eq:detint-orden-base}
 f\prec e\prec b\prec g.
\end{equation}
On exécute l'algorithme interne suivant.
\begin{enumerate}
 \item Chercher les coefficients qui sont des unités de \(\mathbf Z_3\).
 \item Choisir le premier selon \eqref{eq:detint-orden-base} et l'ordre des
 chemins.
 \item Annuler le couple correspondant et enregistrer les opérations
 élémentaires, leur signe et leur contribution à l'orientation.
 \item Répéter jusqu'à ce qu'il ne reste aucun pivot unitaire annulable.
\end{enumerate}
La condition de survie de \eqref{eq:detint-dominio} exige que,
cofinalement en \(m\), le résultat soit un bloc à deux termes
\begin{equation}\label{eq:detint-sx}
 \mathbf Z_3^{r_x}\xrightarrow{S_x}\mathbf Z_3^{r_x},
 \qquad
 \det S_x\ne0\text{ dans }K_3,
\end{equation}
et qu'un raffinement n'ajoute que des blocs contractiles unitaires et des changements
de base enregistrés. Si ces conditions échouent, l'histoire n'appartient pas à
\(\mathcal C_{\rm cont}^{\rm det}\) ; on ne force pas un terminal inexistant.

Le chemin identité fournit un témoin élémentaire. Après le retrait de la racine,
le complexe local est
\begin{equation}\label{eq:detint-testigo-contractil}
 A_nf\xrightarrow{\binom{3c}{1}}
 A_ne\oplus A_nb\xrightarrow{(1,-3c)}A_ng.
\end{equation}
Il est contractile par
\begin{equation}\label{eq:detint-contraccion-local}
 s(g)=e,
 \qquad
 s(ae+\beta b)=\beta f,
 \qquad
 s(f)=0.
\end{equation}
Ainsi, les conditions terminales admettent au moins le témoin identité
interne et ne décrivent pas un sous-domaine formellement vide.

\section{Smith, ordre \texorpdfstring{\(3\)}{3}-adique et unité initiale}

Dans des bases orientées, une diagonalisation par valuation a la forme
\begin{equation}\label{eq:detint-smith}
 U_xS_xV_x
 =\operatorname{diag}(3^{a_1}u_1,\ldots,3^{a_r}u_r),
 \qquad
 a_i\in\mathbf N,quad u_i\in\mathbf Z_3^\times.
\end{equation}
On définit
\begin{equation}\label{eq:detint-orden-d}
 d_x:=v_3(\det S_x)=\sum_{i=1}^r a_i
\end{equation}
et l'unité initiale
\begin{equation}\label{eq:detint-leading-unit}
 \lambda_x:=3^{-d_x}\det S_x\in\mathbf Z_3^\times.
\end{equation}
La coordonnée d'orientation trivialise la droite déterminantielle. Si une
transition produit
\begin{equation}\label{eq:detint-estabilidad-bloques}
 S_x\oplus I_q=U S_{x'}V,
\end{equation}
l'orientation est transportée par le facteur inverse de \(\det U\det V\).
L'unité orientée \(\lambda_x^{\rm or}\) est alors naturelle et ne change pas
lors de l'ajout de blocs contractiles.

En précision finie,
\begin{equation}\label{eq:detint-smith-finito}
 S_{x,n}=S_x\bmod3^n,
 \qquad
 a_i^{(n)}=\min(a_i,n).
\end{equation}
Le terminal est donc la famille compatible de toutes les précisions et
non un entier isolé.

\section{Tour de Bockstein et retenue d'ordre supérieur}

Soit
\[
 C_x^1\xrightarrow{S_x}C_x^0
\]
le bloc terminal. Si une classe \([\bar y]\) admet un relèvement
\(y\in C_x^1\) avec \(S_xy\in3^qC_x^0\), nous définissons
\begin{equation}\label{eq:detint-beta-q}
 \beta_q[\bar y]
 :=\left[3^{-q}S_xy\bmod3\right].
\end{equation}
Le changement de relèvement modifie le représentant par un bord de la page
correspondante, de sorte que \(\beta_q\) est bien défini. La retenue
supérieure est
\begin{equation}\label{eq:detint-carry-superior}
 \operatorname{Carr}_q(y):=3^{-q}S_xy.
\end{equation}
Elle ne remplace pas la retenue TPK de \eqref{eq:detint-c-v-lifts} ; elle enregistre son
prolongement dans la tour des coefficients.

Pour un bloc \(3^{a_i}u_i\), la classe persiste jusqu'à la page \(a_i\),
et à cette page, la différentielle non nulle est la multiplication par
\(\bar u_i\in\mathbf F_3^\times\). Par conséquent,
\begin{equation}\label{eq:detint-bock-longitudes}
 \{\text{longueurs de survie}\}=\{a_1,\ldots,a_r\},
 \qquad
 d_x=\sum_i a_i.
\end{equation}
Le produit orienté des trivialisations de page satisfait
\begin{equation}\label{eq:detint-bock-leading}
 \Theta_{\rm Bock}(x)=\overline{\lambda_x^{\rm or}}
 \in\mathbf F_3^\times.
\end{equation}
La preuve se fait dans chaque bloc diagonal et se transporte par les changements
unimodulaires enregistrés. Smith, Bockstein et l'unité initiale sont ainsi trois
lectures du même terminal.

\iffalse
\section{Les treize coordonnées de la restriction déterminantielle}

Pour chaque \((m,n)\), on définit
\begin{equation}\label{eq:detint-d-trece}
 D_{{\rm det},m,n}(x):=
 (\mathcal F,\operatorname{Ext}_\Gamma,\operatorname{Rel},
 \mathcal N,\operatorname{or},\operatorname{Carr},\operatorname{Mem},
 \partial,\gamma,\operatorname{Msect},\operatorname{Surv},
 \operatorname{Term},\mathcal L)_{{\rm det},m,n}(x),
\end{equation}
avec le contenu suivant.
\begin{enumerate}
 \item La fibre est
 \[
  \mathcal F=(A_n,\mathsf T_m(x),
  \{P_n^0(c_y)\}_y,\{P^{\rm gen}_{y;m,n}\}_y).
 \]
 \item \(\operatorname{Ext}_\Gamma\) contient tous les morphismes
 \(\alpha\), les applications \(\Psi_\alpha\), les homotopies \(K_\alpha\) et
 les applications \(R_{m',m}^{n',n}\).
 \item \(\operatorname{Rel}\) contient \(\partial^2=0\), l'équation de
 pullback, AW, la counité, \(\Psi_\beta\Psi_\alpha=\Psi_{\beta\alpha}\) et
 \(K_{\beta\alpha}=K_\beta+\Psi_\beta K_\alpha\).
 \item
 \(\mathcal N=(m,n,\widetilde c_y,\widetilde v_y,c_{y,n},v_{y,n},
 r_x,a_1,\ldots,a_{r_x},d_x)\).
 \item \(\operatorname{or}\) contient \(y\mapsto\bar y\),
 \(\mathfrak o_G\), \(\Omega\) et l'orientation de la droite
 déterminantielle.
 \item \(\operatorname{Carr}\) conserve séparément la retenue entière,
 sa réduction \(c_{y,n}\) et toutes les retenues supérieures
 \(\operatorname{Carr}_q\).
 \item \(\operatorname{Mem}\) contiene
 \[
  y_0\to\cdots\to y_m,
  \quad\{(\widetilde c_{y_j},\widetilde v_{y_j},\mathbf1_{y_j})\}_{j=0}^m,
  \quad\{K_{\alpha_j}\}.
 \]
 \item
 \(\partial=(\partial_G,\partial_P,\partial_{P^{\rm gen}},
 \partial_{\overline P},h_*)\).
 \item \(\gamma=[y_0\to\cdots\to y_q]\) conserve le simplexe orienté
 complet.
 \item \(\operatorname{Msect}\) contient toutes les extensions
 compatibles, tous les relèvements \(3\)-adiques et toutes les présentations
 unimodulaires du même terminal orienté ; il n'en choisit aucune a posteriori.
 \item \(\operatorname{Surv}\) enregistre la profondeur arbitraire,
 la compatibilité des coefficients, le bilan, l'acyclicité sur \(K_3\),
 la stabilité par blocs unitaires et la permanence de la racine.
 \item
 \[
  \operatorname{Term}=(\widehat P_x^{\rm red},S_x,
  \{a_i,u_i\},d_x,\lambda_x^{\rm or},
  \{\beta_q\},\Theta_{\rm Bock}).
 \]
 \item Le lecteur interne est
 \[
  \mathcal L_{\rm det}^{\rm int}
  =(\mathbf1_y,
  \varpi_{\rm root}z_-=\varpi_{\rm root}z_+,
  [z_-]=[z_+],d_x,\lambda_x^{\rm or},\Theta_{\rm Bock}).
 \]
\end{enumerate}

Les transitions satisfont, coordonnée par coordonnée,
\begin{equation}\label{eq:detint-naturalidad-trece}
 u^{\rm det}_{(m',n'),(m,n)}D_{{\rm det},m',n'}(x')
 =D_{{\rm det},m,n}(\tau_{m',m}x').
\end{equation}

\section{Assembleur, publicateur et chaîne non ablative}

On définit
\begin{equation}\label{eq:detint-graph-d}
 \mathcal G_{\rm det}:=\operatorname{Graph}(D_{\rm det}),
 \qquad
 \operatorname{Res}_{\rm det}(x):=(x,D_{\rm det}x).
\end{equation}
Soit \(\chi_{\rm det}(x)\) la multisection complète de germes
\[
 (y,\widetilde c_y,\widetilde v_y,\operatorname{or}_y,\gamma_y)
\]
compatibles à toute profondeur. L'objet dépendant est
\begin{equation}\label{eq:detint-x-dependiente}
 X_{\chi,{\rm det}}
 :=\{(\chi_{\rm det}(x),x,D_{\rm det}x):
 x\in\mathcal C_{\rm cont}^{\rm det}\},
\end{equation}
et
\begin{equation}\label{eq:detint-prx}
 \operatorname{pr}_{X,{\rm det}}(x,D_{\rm det}x)
 :=(\chi_{\rm det}(x),x,D_{\rm det}x).
\end{equation}

L'assembleur brut
\begin{equation}\label{eq:detint-ensamblador-crudo}
 a_{\rm det}:X_{\chi,{\rm det}}\longrightarrow
 \mathscr M_{\rm det}^{\rm amb}
\end{equation}
produit
\begin{equation}\label{eq:detint-ensamblador-salida}
\begin{split}
 a_{\rm det}(z)=\bigl(&
 \{G_{m,n},\Delta_{\rm AW},\epsilon_G\}_{m,n};\\
 &\{P_n^0(c),\Psi_\alpha,K_\alpha,P^{\rm gen},
 \mathbf1,z_-,z_+,h_*\}_{m,n};\\
 &\{\widehat P^{\rm red},S,d,\lambda^{\rm or},
 \beta_q,\Theta_{\rm Bock}\}\bigr).
\end{split}
\end{equation}
Le macro-objet conserve son entrée :
\begin{equation}\label{eq:detint-graph-a}
 \mathfrak M_{\rm det}:=\operatorname{Graph}(a_{\rm det}),
 \qquad
 \mathcal A_{\rm det}(z):=(z,a_{\rm det}z).
\end{equation}

Le publicateur brut
\begin{equation}\label{eq:detint-publicador-crudo}
 p_{\rm det}:\mathfrak M_{\rm det}\longrightarrow\mathscr Y_{\rm det}
\end{equation}
publie conjointement le certificat
\begin{equation}\label{eq:detint-certificado-publicado}
\begin{gathered}
 \partial^2=0,quad \text{AW et counité},quad
 \mathbf1_y\text{ group-like compatible},\\
 \partial z_\pm=0,quad
 \varpi_{\rm root}z_-=\varpi_{\rm root}z_+,quad
 \partial h_*=z_--z_+,\\
 \Psi_\beta\Psi_\alpha=\Psi_{\beta\alpha},quad
 K_{\beta\alpha}=K_\beta+\Psi_\beta K_\alpha,\\
 d=v_3(\det S)=\sum_i a_i,quad
 \Theta_{\rm Bock}=\overline{\lambda^{\rm or}},quad
 \text{naturalité complète par raffinement}.
\end{gathered}
\end{equation}
Enfin,
\begin{equation}\label{eq:detint-graph-p}
 \mathcal C_{\rm det}^{\rm HMT}:=\operatorname{Graph}(p_{\rm det}),
 \qquad
 \mathcal P_{\rm det}(M):=(M,p_{\rm det}M).
\end{equation}
La chaîne intégrale est
\begin{equation}\label{eq:detint-cadena-completa}
 \mathcal C_{\rm cont}^{\rm disc}\supset
 \mathcal C_{\rm cont}^{\rm det}
 \xrightarrow{\operatorname{Res}_{\rm det}}\mathcal G_{\rm det}
 \xrightarrow{\operatorname{pr}_{X,{\rm det}}}X_{\chi,{\rm det}}
 \xrightarrow{\mathcal A_{\rm det}}\mathfrak M_{\rm det}
 \xrightarrow{\mathcal P_{\rm det}}\mathcal C_{\rm det}^{\rm HMT}.
\end{equation}
Chaque flèche possède un inverse sur son image donné par une projection de
coordonnées. On ne remplace pas l'histoire par une valeur, la multisection par un sélecteur,
le complexe par un déterminant ni la preuve par un nom.

\section{Théorème de réalisation déterminantielle interne}

\begin{theorem}[Réalisation déterminantielle interne complète]
\label{thm:detint-realizacion}
Pour tout \(x\in\mathcal C_{\rm cont}^{\rm det}\), la chaîne
\eqref{eq:detint-cadena-completa} construit un macro-objet naturel en
profondeur et en précision. Il contient deux cycles à racine commune, un filler
explicite qui les identifie, une unité group-like compatible et un terminal
\(3\)-adique dont l'ordre, la forme de Smith et la tour de Bockstein produisent le même
lecteur orienté. Les treize coordonnées de \eqref{eq:detint-d-trece}
restent récupérables en sortie.
\end{theorem}

\begin{proof}
Les équations \eqref{eq:detint-p0-diferencial}--
\eqref{eq:detint-d2} prouvent que \(P_n^0(c)\) est un complexe. Les deux
augmentations sont des applications de complexes, de sorte que le pullback est typé.
Alexander--Whitney est naturel et rend chaque sommet group-like ; avec
\(r\), cela donne l'unité commune. Le calcul direct
\eqref{eq:detint-zpm-ciclos}--\eqref{eq:detint-filler-identidad} prouve que
les publications sont des cycles, ont la même racine et sont reliées par
\(h_*\). Les formules de \(\Psi_\alpha\) prouvent la fonctorialité, et
\(K_\alpha\) prouve la naturalité homotopique en conservant l'accroissement de
mémoire. La définition de \(\mathcal C_{\rm cont}^{\rm det}\) garantit que
la réduction terminale produit un bloc carré sans le postuler hors du
domaine. Smith sur \(\mathbf Z_3\) donne \(d=\sum_i a_i\) ; le calcul de chaque
bloc \(3^{a_i}u_i\) donne
\(\Theta_{\rm Bock}=\overline{\lambda^{\rm or}}\). Troncature, réduction et
ajout de blocs unitaires conservent ces identités. Enfin, les
trois constructions par graphe retiennent l'entrée comme première coordonnée,
de sorte que la composition complète est non ablative.
\end{proof}

\section{Provenance et falsificateurs}

\paragraph{Provenance.}
Le nerf TPK, Alexander--Whitney, la counité, l'unité commune, la forme
normale \(\partial f=3ce+b\), \(\partial e=g\),
\(\partial b=-3cg\), la racine commune, le filler et l'architecture de tour
ont le statut \texttt{ARCHITECTURE\_AUTORALE\_PRÉEXISTANTE}. Le calcul
littéral des deux cycles, de leur racine et de l'identité
\(\partial h_*=z_--z_+\) a le statut \texttt{RÉSULTAT\_RÉCUPÉRÉ}. Le
système local \(\Psi_\alpha/K_\alpha\), le sous-domaine survivant
équilibré, le terminal exclusivement \(3\)-adique et la mise en paquet par
graphes ont le statut \texttt{FORMALISATION\_NOUVELLE} ; ils formalisent sans changer
la provenance conceptuelle de l'appareil.

\begin{ablation}[Falsificateurs internes de la branche déterminantielle]
\label{abl:detint-falsadores}
La construction est falsifiée au niveau correspondant si l'un des
faits suivants se produit.
\begin{enumerate}
 \item \(\partial^2=0\) échoue ; le complexe local n'existe alors pas.
 \item Alexander--Whitney ne commute pas avec le raffinement, la counité échoue ou
 le sommet cesse d'être group-like.
 \item \(z_\pm\) ne sont pas des cycles, leurs racines diffèrent ou
 \(\partial h_*=z_--z_+\) échoue.
 \item \(\Psi_\alpha\) n'est pas une application de complexes ou
 \(K_{\beta\alpha}=K_\beta+\Psi_\beta K_\alpha\) échoue ; dans ce cas, on a
 perdu de la retenue ou de la mémoire.
 \item Le complexe réduit n'est pas équilibré ou n'est pas acyclique sur
 \(K_3\) ; l'histoire reste hors du domaine terminal typé.
 \item Les longueurs de Bockstein ne coïncident pas avec les \(a_i\), ou
 \(\Theta_{\rm Bock}\ne\overline{\lambda^{\rm or}}\).
 \item Le raffinement modifie \(d\) ou \(\lambda^{\rm or}\) en dehors de
 blocs unitaires et de changements d'orientation enregistrés.
 \item L'une des treize coordonnées est omise : \emph{objet substitué ;
 conclusion non transférable}.
\end{enumerate}
Ces falsificateurs vérifient l'objet interne construit ; ils ne le remplacent pas
par un contrôle ultérieur.
\end{ablation}
\fi
